\subsection*{从最后的误差倒推每个中间量的影响}
前向计算从输入走向预测，反向传播则从损失出发，询问每个中间量改变一点会让损失怎样变。输出端已经知道“希望预测往哪里改”，链式法则把这份敏感性逐层传回更早的节点。

沿一条路径，每经过一个运算都乘上它的局部导数；若一个中间量通过多条路径影响损失，要把贡献相加。这个规则来自复合与和的求导法则，并不是神经网络额外发明的数学。

\subsection*{共享参数为什么必须累加贡献}
让同一参数出现在 $L(w)=(w-1)^2+(2w)^2$ 两个分支中。第一条路径贡献 $2(w-1)$，第二条贡献 $8w$。只看第一条会把最优点误判为一；两条相加得到 $10w-2$，真正最优点为 $1/5$。

反向计算的效率来自复用：同一个中间节点到最终损失的敏感性算一次，就能继续传给它的多个输入。若每个参数都从头展开整个损失，会重复大量共同计算。计算图把这些共享结构明确地记录下来。

一次反向传播给出梯度，还没有改变参数；优化器之后才使用梯度更新。若批量损失取平均，梯度也要平均，否则批量大小一变就暗中改变了更新尺度。下一节将把“沿路径乘、汇合时加”写成一个完整的小例子。
